\documentclass[11pt]{article}
\usepackage[a4paper,margin=1in]{geometry}
\usepackage{amsmath,amssymb,amsthm,microtype}
\newtheorem{theorem}{Theorem}[section]
\newtheorem{corollary}[theorem]{Corollary}
\theoremstyle{remark}\newtheorem{remark}[theorem]{Remark}
\title{The Optimal Endpoint Hardy Error under a Spectral Mass Budget}
\author{Bin Wang}\date{July 2026}
\begin{document}\maketitle
\begin{abstract}
We solve the endpoint approximation problem inside the genuine finite normal
spectral class.  If all eigenvalues lie in $|\mu|\le q$ and their squared
mass is at most $M$, the least possible endpoint $H^2$ mismatch has squared
norm asymptotic to $d/\log M$, where $d=\log(q_*/q)$.  The lower bound uses
the all-order deterministic coefficient law; the upper bound uses exact
integer spectral prefixes.  The theorem gives a sharp information-class
constant but does not prove that the actual noisy family converges.
\end{abstract}

\section{Extremal spectral error}
For a finite conjugate multiset $\mathcal S\subset\{|\mu|\le q\}$, put
\[
 \tau_n(\mathcal S)=\sum_{\mu\in\mathcal S}\mu^n,
 \qquad M_2(\mathcal S)=\sum_{\mu\in\mathcal S}|\mu|^2.
\]
Define
\[
 E_{\rm spec}(M)=\inf_{M_2(\mathcal S)\le M}
 \left(\sum_{n\ge2}
 \frac{|\tau_n(\mathcal S)-a_n|^2\rho_*^{2n}}{n^2}
 \right)^{1/2}.
\]
Let $s=q_*/q$ and $d=\log s$.

\begin{theorem}[sharp endpoint mass law]
As $M\to\infty$,
\[
 \boxed{E_{\rm spec}(M)^2\sim\frac{d}{\log M}.}
\]
\end{theorem}
\begin{proof}
For every admissible spectrum, $|\tau_n|\le Mq^{n-2}$.  Fix
$\varepsilon>0$ and put
\[
 n_M=\left\lceil\frac{(1+\varepsilon)\log M}{d}\right\rceil.
\]
Uniformly for $n\ge n_M$,
\[
 \frac{|\tau_n|}{q_*^n}
 \le q^{-2}Ms^{-n}\le q^{-2}M^{-\varepsilon}\longrightarrow0,
\]
whereas RH-268 gives $a_n/q_*^n\to1$.  Therefore
\[
 E_{\rm spec}(M)^2
 \ge(1-o(1))\sum_{n\ge n_M}\frac1{n^2},
\]
and hence
\[
 \liminf_{M\to\infty}\log M\,E_{\rm spec}(M)^2
 \ge\frac d{1+\varepsilon}.
\]
Letting $\varepsilon\downarrow0$ proves the lower bound.

Conversely, RH-316 gives a spectrum matching the first $N$ anchors exactly,
and RH-317 bounds its mass by $Cs^N$.  Its endpoint-scaled spectral tail
satisfies
\[
 \sum_{n>N}\frac{|\tau_n|^2\rho_*^{2n}}{n^2}
 \le C'\sum_{n>N}\frac{s^{2(N-n)}}{n^2}=O(N^{-2}),
\]
whereas the target tail has squared norm
\[
 \sum_{n>N}\frac{|a_n|^2\rho_*^{2n}}{n^2}
 =N^{-1}(1+o(1)).
\]
Thus the mismatch squared norm is $N^{-1}(1+o(1))$.  For a budget $M$,
choose the largest $N$ with $Cs^N\le M$; then
$N=(\log M)/d+O(1)$, proving the upper bound and the asymptotic.
\end{proof}

Numerically,
\[
 d=0.3377217782684642\ldots,
 \qquad \sqrt d=0.5811383469264992\ldots.
\]

\begin{corollary}[actual-family universal lower-rate bound]
If the actual complement has squared mass $M_\sigma\le\sigma^{-1}$, then
\[
 \liminf_{\sigma\downarrow0}
 \log(1/\sigma)\|g_\sigma\|_{H^2(\rho_*)}^2\ge d.
\]
\end{corollary}
\begin{proof}
Repeat the lower-bound argument with
$n_\sigma=\left\lceil
 (1+\varepsilon)\log(1/\sigma)/d\right\rceil$.  The actual mass
cap gives, uniformly for $n\ge n_\sigma$,
\[
 \frac{|\tau_{\sigma,n}|}{q_*^n}
 \le q^{-2}\sigma^{-1}s^{-n}\le q^{-2}\sigma^\varepsilon\to0.
\]
The same reciprocal-square tail yields $d/(1+\varepsilon)$, and then
$\varepsilon\downarrow0$ gives the claim.  This direct proof applies to the
countable Hilbert--Schmidt complement spectrum as well as to finite spectra.
\end{proof}

\begin{remark}
The lower bound tends to zero and is compatible with both convergence and
nonconvergence.  The constructive upper spectra are not identified with the
actual noisy complement.  Gates A--E remain false/open.
\end{remark}
\end{document}
